\section{实验复现与发布}

\subsection{从数据到可部署模型的流水线}

复现实验包含三道关卡：数据收集、训练循环、部署 guard。其中每一关都要持续监控 metrics（KL、entropy、critic loss），并维护实验日志来比较不同 run。

\begin{itemize}
    \item \textbf{数据}：使用 deterministic rollout 保存 state/action/reward/truncation；
    \item \textbf{训练}：搭建 actor/critic/stage scheduler、自动保存 best checkpoint；
    \item \textbf{部署}：将 conservative critic 作为 inference-time guard，保留 fallback policy。
\end{itemize}

\begin{importantbox}{复现流水线的三步守则}
1) 先确保 replay buffer 能被序列化；2) 保持 actor update 与 critic update 在同一 config；3) 在 guard stage 触发 rollback 时记录当时的 KL 和 entropy。
\end{importantbox}

\subsection{版本管理与回滚策略}

为了快速回滚，课程建议用 git tag 记录每个 checkpoint 对应的 config（learning rate、lambda、entropy bonus）。若 new run 出现 high KL，就 revert 到 tag，避免误合并。

\begin{warningbox}{回滚策略的误区}
不要在 rollback 过程中重置 replay buffer，否则会丢失之前的 critic signal；只需回到 checkpoint，保持 buffer state，重新 warm up entropy。
\end{warningbox}

\subsection{本章小结}

实验复现要把训练、日志、回滚打包成一个可执行 pipeline，这样才能在多个团队 copy/paste 同一个好结果。

